\chapter*{Gradient Methods for Model Optimization}
\addcontentsline{toc}{chapter}{Gradient Methods for Model Optimization}

\section*{The Optimization Problem}

Once a model hypothesis $f_{\mathbf{w}}(\xx)$ and a loss function $\ell(\hat{y}, y)$ are chosen, the machine learning task reduces to \textbf{Empirical Risk Minimization (ERM)}: find the parameter vector $\mathbf{w}^*$ that minimizes average loss over the training set:
$$\mathbf{w}^* = \arg\min_{\mathbf{w}} L(\mathbf{w}) = \arg\min_{\mathbf{w}} \frac{1}{N} \sum_{i=1}^N \ell\bigl(f_{\mathbf{w}}(\xx^{(i)}),\; y^{(i)}\bigr)$$

Some linear models (e.g., OLS) admit a closed-form solution via $\nabla_{\mathbf{w}} L = \mathbf{0}$. Most models---logistic regression, neural networks, SVMs with kernels---yield loss surfaces with no analytical solution, requiring \textbf{iterative numerical optimization}.

\section*{Gradient Descent Fundamentals}

\begin{definition}[Gradient Descent (GD)]
Gradient Descent is an iterative first-order optimization algorithm. At each step it evaluates the gradient $\nabla_{\mathbf{w}} L(\mathbf{w})$ at the current point and moves in the direction of steepest descent.
\end{definition}

The update rule at iteration $t$:
$$\mathbf{w}^{(t+1)} = \mathbf{w}^{(t)} - \eta \,\nabla_{\mathbf{w}} L(\mathbf{w}^{(t)})$$

Below is a from-scratch implementation on the convex quadratic $f(x_1,x_2)=\tfrac{1}{2}x_1^2 + \tfrac{5}{2}x_2^2$, whose gradient is $\nabla f = (x_1,\; 5x_2)^T$ and global minimum is at the origin:

\begin{lstlisting}[style=pythonstyle]
import numpy as np

# Define a simple convex quadratic and its gradient
def f(x):
    return 0.5 * x[0]**2 + 2.5 * x[1]**2

def grad_f(x):
    return np.array([x[0], 5 * x[1]])

# Gradient descent from scratch
def gradient_descent(grad, x0, lr=0.1, n_iters=50, tol=1e-6):
    x = x0.copy()
    history = [x.copy()]
    for _ in range(n_iters):
        g = grad(x)
        x = x - lr * g               # GD update rule
        history.append(x.copy())
        if np.linalg.norm(g) < tol:   # stop if gradient ~0
            break
    return x, history

x_opt, path = gradient_descent(grad_f, np.array([2.0, 0.4]))
print(f"Minimum at: {x_opt}")        # -> ~[0, 0]
\end{lstlisting}

\subsection*{The Learning Rate ($\eta$)}

The \textbf{learning rate} (step size) $\eta$ controls how far each update moves along the negative gradient. Even with fixed $\eta$, the actual step length in parameter space is $\|\mathbf{w}^{(t+1)} - \mathbf{w}^{(t)}\| = \eta \cdot \|\nabla L(\mathbf{w}^{(t)})\|$, so steps shrink naturally near a minimum where gradients vanish.

\begin{itemize}
    \item \textbf{Too small} $\eta$: convergence is painfully slow; the optimizer may stall in sub-optimal regions.
    \item \textbf{Too large} $\eta$: updates overshoot the minimum, causing oscillation or divergence.
\end{itemize}

The following snippet demonstrates both failure modes on the same quadratic:

\begin{lstlisting}[style=pythonstyle]
import numpy as np

def grad_f(x):
    return np.array([x[0], 5 * x[1]])

x0 = np.array([2.0, 0.4])

# Too-small learning rate: barely moves in 200 steps
x = x0.copy()
for _ in range(200):
    x = x - 0.001 * grad_f(x)
print(f"lr=0.001 after 200 steps: {x}")   # still far from [0,0]

# Too-large learning rate: diverges!
x = x0.copy()
for _ in range(10):
    x = x - 0.5 * grad_f(x)
print(f"lr=0.5  after 10 steps:  {x}")    # explodes

# Good learning rate: converges smoothly
x = x0.copy()
for _ in range(50):
    x = x - 0.1 * grad_f(x)
print(f"lr=0.1  after 50 steps:  {x}")    # ~[0, 0]
\end{lstlisting}

\begin{remark}[Convexity vs.\ Non-Convexity]
If $L(\mathbf{w})$ is \textbf{convex} (e.g., linear/logistic regression losses), GD converges to the unique global minimum. If $L$ is \textbf{non-convex} (e.g., neural networks), the algorithm may get trapped in \textit{local minima}, \textit{saddle points}, or flat \textit{plateaus}. In practice, SGD noise and adaptive methods help escape shallow traps.
\end{remark}

\section*{Data-Sampling Variants}

Computing $\nabla L$ requires summing over training examples. How many examples per update defines three variants:

\begin{description}
    \item[Batch Gradient Descent (BGD):] Uses the \textbf{entire} dataset ($N$ samples) per update.
    \begin{itemize}
        \item \textit{Pros:} Accurate gradient estimate; stable, predictable convergence.
        \item \textit{Cons:} Computationally expensive per step; impractical for large $N$.
    \end{itemize}

    \item[Stochastic Gradient Descent (SGD):] Uses \textbf{one randomly selected} example per update.
    \begin{itemize}
        \item \textit{Pros:} Extremely fast per iteration; low memory. The noise can help escape shallow local minima.
        \item \textit{Cons:} High variance causes the loss to fluctuate wildly; may never settle precisely at the minimum.
    \end{itemize}

    \item[Mini-Batch Gradient Descent:] The standard compromise---uses a random subset of $B$ examples (e.g., $B \in \{32, 64, 128\}$).
    \begin{itemize}
        \item \textit{Pros:} Reduces update variance vs.\ SGD; exploits GPU vectorization, making it faster than BGD in wall-clock time.
    \end{itemize}
\end{description}

\begin{remark}[Terminology]
In modern deep learning, ``SGD'' almost always refers to mini-batch GD. Pure single-sample SGD is rarely used in practice.
\end{remark}

The conceptual difference in code: BGD computes gradients over all $N$ rows; SGD over 1; mini-batch over a random slice of size $B$:

\begin{lstlisting}[style=pythonstyle]
import numpy as np

# Assume X (N x d), y (N,), w (d,) for linear regression
# MSE gradient: (1/m) * X^T (Xw - y)

def grad_batch(X, y, w):
    """Full-batch gradient (all N samples)."""
    r = X @ w - y
    return X.T @ r / len(y)

def grad_sgd(X, y, w):
    """Stochastic gradient (1 random sample)."""
    i = np.random.randint(len(y))
    r = X[i] @ w - y[i]
    return X[i] * r               # single-sample gradient

def grad_minibatch(X, y, w, B=32):
    """Mini-batch gradient (B random samples)."""
    idx = np.random.choice(len(y), B, replace=False)
    r = X[idx] @ w - y[idx]
    return X[idx].T @ r / B
\end{lstlisting}

\section*{Advanced Gradient Methods}

Standard GD can struggle with pathological curvature (e.g., long narrow ravines) and sensitivity to learning rate choice. Advanced optimizers modify the update rule to accelerate and stabilize convergence.

\subsection*{Momentum-Based Methods}

\begin{definition}[Momentum]
Instead of relying solely on the current gradient, Momentum accumulates an exponentially decaying moving average of past gradients. It acts like a ball rolling downhill---gaining speed in consistent directions and dampening oscillations in erratic ones.
\end{definition}

\begin{description}
    \item[Standard Momentum:] Maintains a velocity vector:
    $$v_t = \gamma \, v_{t-1} + \eta \,\nabla_{\mathbf{w}} L(\mathbf{w}^{(t)})$$
    $$\mathbf{w}^{(t+1)} = \mathbf{w}^{(t)} - v_t$$
    The decay factor $\gamma \in [0,1)$ (typically $0.9$) controls how much past gradients influence the current update.

    \item[Nesterov Accelerated Gradient (NAG):] A ``look-ahead'' variant. Instead of computing the gradient at $\mathbf{w}^{(t)}$, NAG evaluates it at the predicted future position $\mathbf{w}^{(t)} - \gamma \, v_{t-1}$. This anticipatory correction reduces overshooting, especially in regions of steep curvature.
\end{description}

\subsection*{Adaptive Learning Rate Methods}

Rather than one global $\eta$ for all parameters, adaptive methods tune the effective learning rate per parameter $w_j$.

\begin{description}
    \item[AdaGrad:] Scales each parameter's learning rate inversely by the square root of the cumulative sum of its historical squared gradients. Larger updates for infrequent features, smaller for frequent ones. \textit{Drawback:} the accumulated sum grows monotonically, eventually shrinking the learning rate to zero and halting training.

    \item[RMSprop:] Fixes AdaGrad's decay problem by replacing the cumulative sum with an \textbf{exponentially decaying average} of squared gradients, keeping the effective learning rate from vanishing.
\end{description}

\subsection*{Adam (Adaptive Moment Estimation)}

\begin{definition}[Adam Optimizer]
Adam is the most widely used default optimizer in deep learning. It combines \textbf{Momentum} (1st moment---gradient direction) with \textbf{RMSprop} (2nd moment---gradient magnitude scaling).
\end{definition}

The Adam update at iteration $t$ consists of four steps:

\begin{enumerate}
    \item \textbf{1st Moment (Momentum):} Exponentially weighted average of past gradients.
    $$m_t = \beta_1 \, m_{t-1} + (1 - \beta_1)\,\nabla f(\mathbf{w}_{t-1})$$

    \item \textbf{2nd Moment (Adaptive Scaling):} Exponentially weighted average of past squared gradients.
    $$v_t = \beta_2 \, v_{t-1} + (1 - \beta_2)\,\bigl(\nabla f(\mathbf{w}_{t-1})\bigr)^2$$

    \item \textbf{Bias Correction:} Since $m_0 = v_0 = 0$, early estimates are biased toward zero. Correct:
    $$\hat{m}_t = \frac{m_t}{1 - \beta_1^t} \qquad \hat{v}_t = \frac{v_t}{1 - \beta_2^t}$$

    \item \textbf{Parameter Update:}
    $$\mathbf{w}_t = \mathbf{w}_{t-1} - \frac{\eta}{\sqrt{\hat{v}_t} + \epsilon}\;\hat{m}_t$$
    where $\epsilon \approx 10^{-8}$ prevents division by zero.
\end{enumerate}

\begin{remark}[Why Adam Works]
Adam's effective step size is fully adaptive. Large gradient magnitudes shrink the step (preventing divergence); small gradients in plateaus enlarge it (preventing stalling). Combined with momentum's directional smoothing, Adam achieves robust convergence across diverse non-convex landscapes. Default hyperparameters ($\beta_1 = 0.9$, $\beta_2 = 0.999$, $\eta = 10^{-3}$) work well in most settings.
\end{remark}

\section*{Gradient Methods in Practice with \texttt{scikit-learn}}

\texttt{scikit-learn} exposes SGD-based linear models via \texttt{SGDClassifier} and \texttt{SGDRegressor}. These use mini-batch SGD under the hood and accept key hyperparameters: learning rate schedule, regularization, and number of passes over the data.

\begin{lstlisting}[style=pythonstyle]
from sklearn.linear_model import SGDRegressor
from sklearn.preprocessing import StandardScaler
from sklearn.datasets import make_regression

# Generate synthetic regression data
X, y = make_regression(n_samples=500, n_features=5, noise=10,
                       random_state=42)

# SGD is sensitive to feature scale -- always standardize
scaler = StandardScaler()
X_scaled = scaler.fit_transform(X)

# SGDRegressor uses squared loss by default (= linear regression)
sgd_reg = SGDRegressor(
    loss='squared_error',
    learning_rate='invscaling',   # eta = eta0 / t^power_t
    eta0=0.01,                    # initial learning rate
    max_iter=1000,
    tol=1e-4,
    random_state=42
)
sgd_reg.fit(X_scaled, y)
print(f"R^2 = {sgd_reg.score(X_scaled, y):.4f}")
print(f"Learned weights: {sgd_reg.coef_}")
\end{lstlisting}

\begin{lstlisting}[style=pythonstyle]
from sklearn.linear_model import SGDClassifier
from sklearn.datasets import make_classification
from sklearn.preprocessing import StandardScaler

X, y = make_classification(n_samples=1000, n_features=10,
                           random_state=42)
X = StandardScaler().fit_transform(X)

# SGDClassifier with log_loss = logistic regression via SGD
clf = SGDClassifier(
    loss='log_loss',              # logistic regression
    learning_rate='optimal',      # auto-tuned schedule
    max_iter=1000,
    random_state=42
)
clf.fit(X, y)
print(f"Accuracy: {clf.score(X, y):.4f}")
\end{lstlisting}
